\begin{methods}
\section{Materials and Methods}\label{sec:methods}

\subsection{Overview}\label{subsec:overview}

We use HELM sequences mapped into the Uni-Mol latent space as model inputs. The model consists of three core components: a causal context encoder, a context-conditioned denoising network, and a ring-bond predictor. We then optimize the framework for target properties through a diffusion-adapted preference optimization procedure. More details can be found in Figure~\ref{fig:architecture}.

\begin{figure*}[!t]%
\centering
\includegraphics[width=\textwidth]{result_data/fig1_7 (1).png}
\caption{Overview of the PepALD framework.
At each autoregressive step $t$, the Causal Context Encoder transforms the previously generated monomer embeddings $(\mathbf{e}_0, \ldots, \mathbf{e}_{t-1})$ into a context vector $\mathbf{h}_t$.
A context-conditioned Diffusion Engine then samples a continuous embedding $\mathbf{z}_t$ via iterative denoising in the Uni-Mol latent space.
The Token Mapper discretizes $\mathbf{z}_t$ into a HELM monomer token $x_t$ subject to position-dependent R-group chemical constraints.
Concurrently, the Autoregressive Ring Predictor evaluates potential ring-closure bonds between $x_t$ and all preceding residues using R-group compatibility scoring.
The process repeats until the target sequence length is reached, yielding a complete macrocyclic peptide with predicted ring topology.
}\label{fig:architecture}
\end{figure*}
%%
%% [DRAWING GUIDE FOR FIGURE 1 — Overall Architecture]
%%
%% Layout: horizontal left-to-right pipeline, with a vertical time axis showing
%% autoregressive unrolling for steps t=0, 1, ..., T-1.
%%
%% LEFT BLOCK — "Causal Context Encoder":
%%   - Show a stack of causal Transformer layers with a triangular attention mask.
%%   - Input: sequence of Uni-Mol monomer embeddings (e_0, ..., e_{t-1}).
%%   - A special "Start Token" embedding feeds position t=0.
%%   - Output arrow labeled h_t pointing right.
%%
%% CENTER BLOCK — "Diffusion Engine":
%%   - Show the reverse diffusion trajectory: z_K (noise) → z_{K-1} → ... → z_0 (clean).
%%   - Inside each denoising step, show a DenoiserBlock with:
%%       (a) AdaptiveLayerNorm conditioned on timestep k,
%%       (b) Cross-Attention receiving h_t from the left block.
%%   - Output arrow labeled z_t pointing right.
%%
%% RIGHT BLOCK — "Token Mapper + Ring Predictor":
%%   - Token Mapper: show z_t compared against a codebook of Uni-Mol reference
%%     embeddings via cosine distance, with a filter icon for R-group constraints.
%%   - Below or beside it, show the Ring Predictor: current R-group embeddings
%%     (R1, R2, R3) dot-producted against history R-groups, producing a 3x3
%%     compatibility matrix, feeding into position and type prediction heads.
%%   - Output: discrete token x_t and optional ring bond annotation.
%%
%% BOTTOM — "Feedback loop":
%%   - Arrow from x_t back to the left block, showing the embedding of x_t
%%     being appended to the history for the next step.
%%
%% OPTIONAL INSET — "Training Objectives":
%%   - Small box showing the three loss terms: L_diff (MSE), L_CE, L_ring.
%%




\subsection{Representing HELM Grammar with Chemically Informed Monomer Embeddings}\label{subsec:helm}

We adopt HELM as the native sequence representation \citep{ref_helm}. Under the HELM specification, each monomer exposes a set of canonical attachment points---such as $\mathrm{R1}$, $\mathrm{R2}$, and $\mathrm{R3}$---through which all covalent bonds, backbone amides and intramolecular ring closures alike, are explicitly specified.

We next associate each monomer in the HELM vocabulary with a pre-computed embedding derived from Uni-Mol, a molecular representation model pre-trained to capture three-dimensional molecular geometry and chemical semantics \citep{ref_unimol}. To preserve both residue-level identity and attachment-site chemistry, we extract a molecule-level representation for the whole monomer and atom-level representations for the atoms adjacent to the HELM R-group attachment sites. Missing R-group sites are represented by zero vectors.
Specifically, for each monomer $m$ with SMILES representation, we extract a structured embedding vector:
\begin{equation}
\mathbf{e}_m = \bigl[\mathbf{e}_m^{\mathrm{CLS}},\; \mathbf{e}_m^{\mathrm{R1}},\; \mathbf{e}_m^{\mathrm{R2}},\; \mathbf{e}_m^{\mathrm{R3}}\bigr] \in \mathbb{R}^{4 \times d}
\label{eq:unimol_full}
\end{equation}
where $\mathbf{e}_m^{\mathrm{CLS}} \in \mathbb{R}^d$ is the molecule-level representation, and $\mathbf{e}_m^{\mathrm{R}i} \in \mathbb{R}^d$ ($i \in \{1,2,3\}$) are the atomic representations at the respective R-group attachment sites.

The embedding matrix is pre-computed for all $|\mathcal{V}|$ monomers in the vocabulary and frozen during training, ensuring that the diffusion process operates in a chemically grounded latent space rather than a learned but potentially arbitrary token space.


\subsection{ALD generative model}\label{subsec:ald_generative_model}

Conditioned on a target length $T$, PepALD defines the probability of a HELM monomer sequence $\mathbf{x}=(x_1,\ldots,x_T)$ by an ALD process. At position $t$, the already generated prefix is summarized as a causal context vector $\mathbf{h}_t=f_{\mathrm{ctx}}(x_{<t})$, with a learnable start context used for $t=1$. The clean latent variable for the next residue is identified with the frozen Uni-Mol embedding of the corresponding monomer, $\mathbf{z}_{t,0}=\tilde{\mathbf{e}}_{x_t}$. Thus,
\begin{equation}
p_\theta(\mathbf{x}\mid T)=\prod_{t=1}^{T}p_\theta(x_t\mid x_{<t}),
\label{eq:ald_ar_factorization}
\end{equation}
where each discrete conditional distribution is induced by a continuous conditional density $p_\theta(\mathbf{z}_{t,0}\mid \mathbf{h}_t)$ followed by a chemically constrained projection to the HELM vocabulary:
\begin{equation}
p_\theta(x_t=m\mid x_{<t})=
\int \mathds{1}\!\left[\mathcal{M}_t(\mathbf{z},\mathbf{h}_t)=m\right]\,
p_\theta(\mathbf{z}\mid \mathbf{h}_t)\,d\mathbf{z}.
\label{eq:ald_induced_token_distribution}
\end{equation}
Here $\mathcal{M}_t$ denotes the position-aware token mapping rule, and its admissible set is restricted by the R-group requirements of the current backbone position.

\subsubsection{Forward diffusion}

For each valid residue position, the forward process perturbs the clean monomer embedding independently in the Uni-Mol latent space. Let $\{\beta_k\}_{k=1}^{K}$ be the variance schedule, $\alpha_k=1-\beta_k$, and $\bar{\alpha}_k=\prod_{s=1}^{k}\alpha_s$. PepALD uses the standard DDPM Gaussian kernel \citep{ref_ddpm}:
\begin{equation}
q(\mathbf{z}_{t,k}\mid \mathbf{z}_{t,0},\mathbf{h}_t)
=\mathcal{N}\!\left(\sqrt{\bar{\alpha}_k}\mathbf{z}_{t,0},
(1-\bar{\alpha}_k)\mathbf{I}\right),
\label{eq:ald_forward_kernel}
\end{equation}
In the forward process, the noising kernel is conditionally independent of the autoregressive context, i.e., $q(\mathbf{z}_{t,k}\mid \mathbf{z}_{t,0},\mathbf{h}_t)=q(\mathbf{z}_{t,k}\mid \mathbf{z}_{t,0})$.
equivalently,
\begin{equation}
\mathbf{z}_{t,k}=\sqrt{\bar{\alpha}_k}\mathbf{z}_{t,0}
+\sqrt{1-\bar{\alpha}_k}\boldsymbol{\epsilon},\quad
\boldsymbol{\epsilon}\sim\mathcal{N}(\mathbf{0},\mathbf{I}).
\label{eq:ald_forward_sample}
\end{equation}
During training, $k$ is sampled uniformly from $\{1,\ldots,K\}$ for every non-padding token, and the denoising network is trained to predict the injected noise $\boldsymbol{\epsilon}$ from $(\mathbf{z}_{t,k},k,\mathbf{h}_t)$.

\subsubsection{Reverse denoising with hybrid sampling}\label{subsubsec:reverse_hybrid_sampling}

Generation starts from isotropic Gaussian noise $\mathbf{z}_{t,K}\sim\mathcal{N}(\mathbf{0},\mathbf{I})$ for each autoregressive position. The reverse transition is parameterized by the context-conditioned noise predictor $\boldsymbol{\epsilon}_\theta$:
\begin{equation}
p_\theta(\mathbf{z}_{t,k-1}\mid \mathbf{z}_{t,k},\mathbf{h}_t)
=\mathcal{N}\!\left(\boldsymbol{\mu}_\theta(\mathbf{z}_{t,k},k,\mathbf{h}_t),
\sigma_k^2\mathbf{I}\right),
\label{eq:ald_reverse_kernel}
\end{equation}
with
\begin{equation}
\boldsymbol{\mu}_\theta=
\frac{1}{\sqrt{\alpha_k}}\left(
\mathbf{z}_{t,k}
-\frac{1-\alpha_k}{\sqrt{1-\bar{\alpha}_k}}\,
\boldsymbol{\epsilon}_\theta(\mathbf{z}_{t,k},k,\mathbf{h}_t)
\right),
\label{eq:ald_reverse_mean}
\end{equation}
and $\sigma_k^2=\beta_k(1-\bar{\alpha}_{k-1})/(1-\bar{\alpha}_k)$, with the stochastic term omitted at $k=1$ \citep{ref_ddpm}. For faster inference, the same predicted noise can be used in a DDIM trajectory over a reduced and adjustable subset of timesteps \citep{ref_ddim}.

After the final denoising step, the continuous sample is projected back to a valid HELM monomer using a hybrid decoder. At the core of the decoder is a fused score that reconciles the geometric output of the diffusion process with the contextual prior of the autoregressive state. Let $\tilde{\mathbf{e}}_m$ denote the $\ell_2$-normalized reference embedding of monomer $m$, and let $\mathbf{h}_t$ denote the causal context representation at position $t$. For every admissible candidate $m \in \mathcal{C}_t$, we define:
\begin{equation}
\mathcal{S}(m \mid \mathbf{z}_{t,0}, \mathbf{h}_t)
= d(\mathbf{z}_{t,0}, \tilde{\mathbf{e}}_m)
- \lambda \log p_{\mathrm{LM}}(m \mid \mathbf{h}_t),
\label{eq:ald_hybrid_sampling}
\end{equation}
where $d(\cdot,\cdot)$ is the cosine distance in the shared embedding space, $p_{\mathrm{LM}}(m \mid \mathbf{h}_t)=\mathrm{softmax}(\mathbf{W}_{\mathrm{LM}}\mathbf{h}_t)_m$ is the categorical distribution produced by an auxiliary language-model head trained jointly with the denoiser, and the fusion weight $\lambda \geq 0$ is manually adjusted during sequence generation.
Related experiments are provided in Section~3.


\subsection{Model architecture}\label{subsec:ald}

As shown in Fig.~\ref{fig:architecture}, PepALD implements the ALD formulation in Section~\ref{subsec:ald_generative_model} with three coordinated modules: context encoder, denoising network and ring-bond predictor.

\subsubsection{Causal context encoder}\label{subsubsec:context_encoder}

To encode the autoregressive history, we use a causal context encoder $f_{\mathrm{ctx}}$.
Given the Uni-Mol embeddings of the generated prefix $(\tilde{\mathbf{e}}_{x_1}, \ldots, \tilde{\mathbf{e}}_{x_{t-1}})$, the encoder outputs a position-specific context vector $\mathbf{h}_t$ that conditions the diffusion model for the next monomer.

The encoder takes the prefix Uni-Mol embeddings with added sinusoidal positional encodings as input.
For efficient teacher-forced training, the embedded HELM sequence is processed in parallel by $L_{\mathrm{ctx}}$ causal Transformer layers \citep{ref_transformer}.
Within this stack, each causal Transformer layer consists of masked multi-head self-attention, a feed-forward block, residual connections, layer normalization, and dropout.
The attention mask is lower triangular, and a padding mask is applied when minibatches contain peptides of different lengths.
Thus, the hidden state at position $i$ can depend only on residues $x_{\leq i}$ that are present in the sequence.

The output of the stack is then passed through an additional final layer normalization and a linear context projection, giving the overall context encoder flow:
\begin{equation}
\mathbf{H} = \mathrm{Proj}_{\mathrm{ctx}}\!\left(\mathrm{LN}\!\left[\mathrm{CausalTransformer}\!\left(\mathrm{PosEnc}(\mathbf{W}_{\mathrm{in}}\tilde{\mathbf{E}})\right)\right]\right).
\label{eq:causal_encoder}
\end{equation}
Here $\tilde{\mathbf{E}}$ denotes the matrix of input Uni-Mol monomer embeddings, and $\mathbf{W}_{\mathrm{in}}$ is the input projection to the Transformer hidden dimension.
For the first residue, where no prefix is available, we introduce a learnable start embedding and pass it through the same encoding pathway to obtain $\mathbf{h}_{\mathrm{start}}$.
For subsequent positions, the context used to predict residue $x_t$ is the previous contextual state $\mathbf{H}_{t-1}$.
In practice, the shifted context sequence is assembled as $[\mathbf{h}_{\mathrm{start}}, \mathbf{H}_1,\ldots,\mathbf{H}_{T-1}]$, yielding one context vector for every valid residue position while preserving the same autoregressive information flow used during generation.

\subsubsection{Context-conditioned denoising network}

Given the context vector $\mathbf{h}_t$, we introduce a context-conditioned denoiser that guides the reverse diffusion trajectory for the next monomer.
The denoiser is designed around two conditioning axes: a normalized timestep signal that marks the current denoising stage, and a context signal that anchors each refinement step to the prefix chemistry.

The noisy embedding $\mathbf{z}_{t,k}$ is first projected to the model dimension and combined with a sinusoidal timestep embedding $\mathbf{e}_k=\mathrm{MLP}_{t}(\mathrm{SinEmb}(k/K))$.
The resulting state is processed by $L_{\mathrm{den}}$ denoiser blocks:
\begin{equation}
\begin{aligned}
\mathbf{u}^{(0)} &= \mathbf{W}_{z}\mathbf{z}_{t,k} + \mathbf{e}_k,\\
\tilde{\mathbf{u}}^{(l)} &= \mathrm{AdaLN}_{l}\!\left(\mathbf{u}^{(l-1)},\, \mathbf{e}_k\right),\\
\mathbf{u}^{(l)} &= \mathrm{DenoiserBlock}_{l}\!\left(\tilde{\mathbf{u}}^{(l)},\, \mathbf{h}_t\right),\quad l=1,\ldots,L_{\mathrm{den}}.
\end{aligned}
\label{eq:denoiser_block}
\end{equation}
Each denoiser block applies a self-attention-based update to the single noisy-token state, cross-attention to $\mathbf{h}_t$, and a feed-forward transformation with residual connections, layer normalization, and dropout.
AdaLN supplies timestep-dependent scale and shift parameters, whereas cross-attention provides the only path through which the autoregressive history enters the denoising network.

The final block output is normalized and projected back to the Uni-Mol embedding dimension to predict the injected noise:
\begin{equation}
\hat{\boldsymbol{\epsilon}}
= \mathbf{W}_{\mathrm{out}}\mathrm{LN}\!\left(\mathbf{u}^{(L_{\mathrm{den}})}\right)
= \boldsymbol{\epsilon}_\theta(\mathbf{z}_{t,k},\, k,\, \mathbf{h}_t).
\label{eq:noise_prediction}
\end{equation}
The resulting noise estimate is used by the DDPM or DDIM samplers described in Section~\ref{subsubsec:reverse_hybrid_sampling}.


\subsubsection{R-group-aware ring bond predictor}\label{subsec:ring}

To couple macrocyclization with residue generation, PepALD predicts ring-closure bonds within the autoregressive loop.
At each step $t$, the predictor scores whether the newly generated residue $x_t$ should connect to a previous residue $x_j$ ($j<t$) and classifies the bond into one of four R-group link types, $\mathcal{B}=\{\mathrm{R3R3},\mathrm{R1R2},\mathrm{R1R3},\mathrm{R3R2}\}$.

For each candidate pair $(x_j,x_t)$, the model combines a sequence-level context feature with an R-group compatibility feature.
The context feature reuses the shifted states $\mathbf{h}_j$ and $\mathbf{h}_t$ produced by the causal context encoder in Section~\ref{subsubsec:context_encoder}, while the R-group feature is computed from the scaled dot products between the Uni-Mol attachment-site embeddings of the current and historical residues.
At a high level, the predictor can be written as
\begin{equation}
\left(s_{j,t}^{\mathrm{pos}},\mathbf{l}_{j,t}^{\mathrm{type}}\right)
=
f_{\mathrm{ring}}\!\left(
\mathbf{h}_t,\mathbf{h}_j,
\mathbf{E}_{x_t}^{\mathrm{R}},\mathbf{E}_{x_j}^{\mathrm{R}}
\right),
\quad j<t,
\label{eq:ring_predictor_summary}
\end{equation}
where $s_{j,t}^{\mathrm{pos}}$ is the bond-position score and $\mathbf{l}_{j,t}^{\mathrm{type}}$ contains the logits over R-group link types.
Full definitions of the pair features and prediction heads are provided in Supplementary Section~6.


\subsection{Supervised pre-training and fine-tuning}\label{subsec:training}

PepALD is trained with supervised objectives that target complementary aspects of peptide sequence generation.
The first stage pre-trains the context encoder and denoiser using denoising and auxiliary token-prediction losses.
The second stage fine-tunes the same generator on curated macrocyclic peptides and adds an autoregressive ring-bond objective for learning cyclization topology.

For a mini-batch, let $\Omega$ denote the set of valid non-padding residue positions and let $\mathbf{z}_{b,t,0}=\tilde{\mathbf{e}}_{x_{b,t}}$ be the clean Uni-Mol embedding of token $x_{b,t}$.
For each $(b,t)\in\Omega$, a diffusion step $k$ and Gaussian noise $\boldsymbol{\epsilon}_{b,t}$ are sampled as in Section~\ref{subsec:ald_generative_model}.
The denoiser is optimized by the standard $\epsilon$-prediction objective \citep{ref_ddpm}:
\begin{equation}
\mathcal{L}_{\mathrm{diff}}
=\frac{1}{|\Omega|d}\sum_{(b,t)\in\Omega}
\left\|
\boldsymbol{\epsilon}_{\theta}(\mathbf{z}_{b,t,k}, k, \mathbf{h}_{b,t})
-\boldsymbol{\epsilon}_{b,t}
\right\|_2^2 ,
\label{eq:diff_loss}
\end{equation}
where $d$ is the Uni-Mol embedding dimension and $\mathbf{h}_{b,t}$ is the shifted context vector used to predict position $t$.
In parallel, the language-model head receives the same context and predicts the ground-truth monomer identity:
\begin{equation}
\begin{aligned}
p_{\mathrm{LM}}(m\mid \mathbf{h}_{b,t})
&=\left[\operatorname{softmax}(\mathbf{W}_{\mathrm{LM}}\mathbf{h}_{b,t})\right]_m,\\
\mathcal{L}_{\mathrm{CE}}
&=-\frac{1}{|\Omega|}\sum_{(b,t)\in\Omega}
\log p_{\mathrm{LM}}(x_{b,t}\mid \mathbf{h}_{b,t}).
\end{aligned}
\label{eq:ce_loss}
\end{equation}
This auxiliary term regularizes the contextual representation used by the hybrid decoder in Section~\ref{subsubsec:reverse_hybrid_sampling}.

During macrocyclic fine-tuning, the ring predictor is trained on all valid autoregressive candidate pairs.
Let $\mathcal{P}=\{(b,j,t): (b,t)\in\Omega,\; j<t\}$ and let $y_{b,j,t}^{\mathrm{pos}}\in\{0,1\}$ indicate whether residue $x_{b,t}$ forms a ring bond with a previous residue $x_{b,j}$.
The bond-position objective is a weighted binary cross-entropy over the raw logits $s_{b,j,t}^{\mathrm{pos}}$:
\begin{equation}
\begin{aligned}
\pi_{b,j,t} &= \sigma(s_{b,j,t}^{\mathrm{pos}}),\\
\mathcal{L}_{\mathrm{pos}}
&=-\frac{1}{|\mathcal{P}|}
\sum_{(b,j,t)\in\mathcal{P}}
\Big[
w^{+}y_{b,j,t}^{\mathrm{pos}}\log \pi_{b,j,t}\\
&\qquad\qquad\qquad
+(1-y_{b,j,t}^{\mathrm{pos}})\log (1-\pi_{b,j,t})
\Big],
\end{aligned}
\label{eq:pos_loss}
\end{equation}
where $w^{+}=\min(n^{-}/n^{+},50)$ is computed from the positive and negative candidate counts in the mini-batch to compensate for sparse ring closures.
For positive candidate pairs $\mathcal{P}^{+}=\{(b,j,t)\in\mathcal{P}:y_{b,j,t}^{\mathrm{pos}}=1\}$, a separate cross-entropy term classifies the R-group link type:
\begin{equation}
\mathcal{L}_{\mathrm{type}}
=-\frac{1}{|\mathcal{P}^{+}|}
\sum_{(b,j,t)\in\mathcal{P}^{+}}
\log
\frac{\exp(l_{b,j,t,c^{*}}^{\mathrm{type}})}
{\sum_{c\in\mathcal{B}}\exp(l_{b,j,t,c}^{\mathrm{type}})} ,
\label{eq:type_loss}
\end{equation}
where $c^{*}$ is the ground-truth bond type and $\mathcal{B}$ is the four-class R-group bond vocabulary defined in Section~\ref{subsec:ring}.
If a mini-batch contains no positive ring pair, the type term is set to zero.

The supervised objective combines the denoising, token-prediction, and ring-bond terms:
\begin{equation}
\mathcal{L}_{\mathrm{sup}}
=\mathcal{L}_{\mathrm{diff}}
+\lambda_{\mathrm{CE}}\mathcal{L}_{\mathrm{CE}}
+\lambda_{\mathrm{ring}}\left(\mathcal{L}_{\mathrm{pos}}+\mathcal{L}_{\mathrm{type}}\right).
\label{eq:sup_loss}
\end{equation}
During pre-training, the ring objective is disabled; during macrocyclic fine-tuning, all three supervised components are active with configurable weights.

\subsection{Optimization for cell-penetrating peptide design}\label{subsec:preference_optimization}

The downstream task is to optimize PepALD for cyclic cell-penetrating peptide design.
We first construct a permeability-enriched prior generator and then align it with a reward-ranked preference objective.

\subsubsection{Permeability prior fine-tuning}\label{subsubsec:permeability_prior}

For permeability-oriented prior construction, we rank the merged ChEMBL32--CycPeptMPDB prior table by predicted permeability and select the top 1,000 entries for an additional cyclic-only fine-tuning stage.
The resulting permeability-enriched prior model serves as the initialization for the trainable preference-aligned policy.

\subsubsection{Composite reward and WP-DPO objective}\label{subsubsec:dpo_loss}

In the target-specific DPO experiments, candidate rewards are computed from Uni-Dock/Vina docking scores \citep{ref_unidock}.
Because lower docking scores indicate stronger predicted binding, the docking-derived reward is normalized by a robust transformation over the valid candidate set before preference-pair construction.
The resulting scalar score is denoted $R(\mathbf{x})$.
Further details on the sign convention, robust normalization, and invalid-score handling are provided in Supplementary Sections~4 and~7.

Direct preference optimization (DPO) has been used to align generative models from pairwise preferences \citep{ref_dpo}, and its diffusion variant replaces sequence likelihoods with denoising errors along the diffusion process \citep{ref_diffusion_dpo}.
Given a preference pair $(\mathbf{x}^{w},\mathbf{x}^{l})$ with $R(\mathbf{x}^{w})>R(\mathbf{x}^{l})$, we sample a diffusion step $k$ and noise $\boldsymbol{\epsilon}$.
For a model $\theta$, $\overline{\mathrm{MSE}}_{\theta}(\mathbf{x})$ denotes the mean squared difference between the predicted and injected noise, averaged over all valid residue positions of peptide $\mathbf{x}$.
Let
\begin{equation}
\Delta_{\theta}(\mathbf{x})
=\overline{\mathrm{MSE}}_{\mathrm{ref}}(\mathbf{x})
-\overline{\mathrm{MSE}}_{\theta}(\mathbf{x})
\label{eq:dpo_delta}
\end{equation}
be the improvement of the current model over the reference model.
The preference loss is
\begin{equation}
\mathcal{L}_{\mathrm{DPO}}
=-\frac{1}{B}\sum_{b=1}^{B}
\log\sigma\!\left(
\beta_{\mathrm{DPO}}
\left[
\Delta_{\theta}(\mathbf{x}_{b}^{w})
-\Delta_{\theta}(\mathbf{x}_{b}^{l})
\right]\right),
\label{eq:dpo_loss}
\end{equation}
where $B$ is the number of preference pairs in a mini-batch and $\beta_{\mathrm{DPO}}$ controls the strength of the preference margin.
For each pair, the winner and loser are evaluated with the same sampled $k$ and $\boldsymbol{\epsilon}$.
Motivated by DPO-Positive (DPOP), which addresses degradation of preferred samples under standard DPO \citep{ref_dpop}, we add an external winner-protection regularizer to prevent the policy from improving the relative margin by degrading the denoising quality of winner samples:
\begin{equation}
\mathcal{L}_{\mathrm{win}}
=\frac{1}{B}\sum_{b=1}^{B}
\max\left(
0,\,
\overline{\mathrm{MSE}}_{\theta}(\mathbf{x}_{b}^{w})
-
\overline{\mathrm{MSE}}_{\mathrm{ref}}(\mathbf{x}_{b}^{w})
\right).
\label{eq:winner_reg}
\end{equation}
The final WP-DPO objective is
\begin{equation}
\mathcal{L}_{\mathrm{align}}
=\mathcal{L}_{\mathrm{DPO}}
+\alpha_{\mathrm{win}}\mathcal{L}_{\mathrm{win}},
\label{eq:dpo_total_loss}
\end{equation}
where $\alpha_{\mathrm{win}}$ controls the strength of winner protection.

\subsubsection{Iterative candidate-pool construction and preference pairing}\label{subsubsec:pairing}

Preference construction is performed iteratively on model-generated candidate pools rather than on the original training corpus.
At each WP-DPO round, candidates are sampled from the preceding PepALD checkpoint, post-processed to enforce the target cyclic topology, deduplicated, filtered for valid Uni-Dock/Vina scoring, and ranked by the target-specific reward described above.
Preference pairs are then constructed with a stage-dependent, diversity-aware strategy.
In early WP-DPO rounds, high- and low-reward regions define winner and loser candidate pools, and diversity-aware selection is applied within each pool before pairing.
In later refinement rounds, after the generator has begun to produce competitive binders, we make the preference task more stringent by replacing obvious low-quality losers with harder negatives.
Specifically, loser candidates are sampled from an intermediate Vina-score range, so that they remain plausible binders but are still worse than the selected winners.
We then apply diversity-aware reranking to both winner and loser pools using cyclic-residue bigram features, and pair each winner with a structurally similar loser \citep{ref_mmr,ref_drugblip}.
Pairs are retained only when the loser remains meaningfully worse than its matched winner according to the combined reward.
The final $(\mathbf{x}^{w},\mathbf{x}^{l})$ pairs are fixed before training and reused unchanged during WP-DPO optimization.
Details of diversity selection, hard-negative pairing, and candidate-pool settings are provided in Supplementary Sections~5 and~9.

\subsection{Datasets}\label{subsec:datasets}

We built the training corpus from public peptide resources that provide residue-level monomer annotations and explicit inter-monomer connections.
The main sequence source was ChEMBL32 \citep{ref_chembl}, from which 20,783 HELM sequences were extracted for pre-training.

Then, we performed macrocyclic peptide fine-tuning using \linebreak[4]CycPeptMPDB \citep{ref_cycpeptmpdb}, a curated database of cyclic peptides with membrane-permeability measurements and structural representations.
The raw CycPeptMPDB HELM strings were normalized to a single polymer identifier, deduplicated, and written to the processed HELM corpus, which yielded 7,348 unique cyclic peptide HELM sequences.
After merging and preprocessing, the training corpus contains predominantly head-to-tail cycles, together with terminal-to-internal cyclization patterns.

The monomer vocabulary was compiled from the curated HELM monomer library used in this study, covering monomers observed in ChEMBL32 and CycPeptMPDB together with additional peptide building blocks.
The library contains 3,105 peptide monomer tokens with available R-group attachment sites.
For every monomer, a Uni-Mol representation was precomputed from its molecular structure, including a molecule-level vector and attachment-site vectors for the available R-groups.

\end{methods}
